\documentclass[10pt,a4paper]{amsart}
\usepackage[margin=19mm]{geometry}
\usepackage{amsmath,amssymb,mathtools}
\usepackage[scr=rsfs,cal=euler]{mathalfa}
\usepackage{xfrac}
\usepackage[unicode,hidelinks,bookmarks=false]{hyperref}
\newcommand{\ie}{i.e.\ }
\newcommand{\cf}{cf.\ }
\newcommand{\resp}{respectively\ }
\newcommand{\Subsection}{\subsection}
\providecommand{\nobreakdash}{\nobreak\hbox{-}}
\setlength{\emergencystretch}{2em}
\multlinegap0pt
\usepackage{amssymb}
\usepackage{mathtools}
\usepackage{tcolorbox}
\usepackage{float}
\usepackage{xcolor}
\usepackage{tikz}
\newtheorem{lemma}{Lemma}
\newtheorem{proposition}{Proposition}
\newcommand \R {{ \mathbb R}}
\newcommand \T {{ \mathbb T}}
\newcommand \cC {{\mathscr{C}}}
\newcommand \cI {{\mathcal I}}
\newcommand*{\diff}{\mathop{}\!\mathrm{d}}
\newcommand{\geo}{\mathsf{g}}
\newcommand{\horo}{\mathsf{h}}
\title{Horocycle flows on abelian covers of surfaces of negative curvature}
\author{Roberto Castorrini, Davide Ravotti}
\date{Source: arXiv:2405.08592v3 (CC BY 4.0)}
\begin{document}
\maketitle

\noindent\emph{Source and licence.} Horocycle flows on abelian covers of surfaces of negative curvature. Version arXiv:2405.08592v3 (CC BY 4.0), distributed by arXiv under the Creative Commons Attribution 4.0 International licence, http://creativecommons.org/licenses/by/4.0/, as recorded in the arXiv licence metadata for this exact version and shown on the abstract page https://arxiv.org/abs/2405.08592v3. Full licence notice: \texttt{licenses/arxiv-2405.08592-CC-BY-4.0.txt}.\par\smallskip

\noindent\textit{Editorial scope.} Counted unit: the article's proposition prop:G (source0.tex lines 2325--2342). The article's complete original proof of it is delivered with it (source0.tex lines 2344--2459, one \texttt{proof} environment of 116 lines). No mathematics is written, repaired or re-derived: the statement and the proof are the article's own lines, in the article's own notation, and no retained line is substituted or re-flowed.  Retained uncounted as the source-local context the proof needs: lines 438--451; lines 448--450; lines 504--510; lines 513--518.  Nothing was omitted from any retained span, so the package carries no omission record.  No retained line contains a citation, so the wrapper carries no bibliography and no bibliography entry.  No retained line contains a figure environment, an image-inclusion command or drawing code, so no image is delivered and the package declares no asset dependency.  Editorial formatting only, in addition: an AMS \texttt{amsart} preamble that restates the article's own declarations and macros used by the retained lines, namely the environments \texttt{lemma}, \texttt{proposition} on separate counters, together with the standard packages those lines need.  The retained environments print in this document as Lemma 1, Proposition 1 in order of appearance, whereas the article prints the same items with the numbering of its own full document.  The complete native source of the article as distributed in the arXiv e-print for this exact version is bundled unchanged as \texttt{sources/159133/source0.tex}.
\begin{lemma}\label{lemma:Jacobi}
For each $x\in M$, the function
\[
J(t,x) := \frac{\partial}{\partial s}\Big\vert_{s=0} \tau(s,t,x)
\]
satisfies
\[
\partial^2_{t}{J}(t,x) + K(\geo_t(x))J(t,x) =0, \qquad J(0,x)=1, \qquad \lim_{t \to \infty} J(t,x) = 0,
\]
where $K$ denotes the curvature. Moreover, given $\overline k, \underline k<0$ such that $\underline k \leq K(\geo_t(x)) \leq \overline k$ for each $x\in M$ and $t\in \R^+$, then
\begin{equation}\label{eq:Jbound}
e^{-\sqrt{-\underline k}t} \le J(t,x) \le e^{-\sqrt{-\overline k}t}, \qquad \forall t\in \R_{>0}, \quad \forall x\in M.
\end{equation}
\end{lemma}

\begin{equation}
e^{-\sqrt{-\underline k}t} \le J(t,x) \le e^{-\sqrt{-\overline k}t}, \qquad \forall t\in \R_{>0}, \quad \forall x\in M.
\end{equation}

 \begin{equation}\label{eq:potential}
 \begin{split}
 J_{-t}(x)&=\exp\left({\int_0^t -\operatorname{div} (X|_{E^-})\circ \geo_{-s}(x)\mathrm d s}\right)=D\geo_{-t}(x)|_{E_-},\\
 %
  J^u_{-t}(x)&=\exp\left({\int_0^{-t} \operatorname{div} (X|_{E^+})\circ \geo_s(x)\mathrm d s}\right)=D\geo_{t}(x)|_{E_+},
 \end{split}
 \end{equation}

\begin{equation}\label{eq:div}
\begin{split}
&\frac{\diff}{\diff t}\Big\vert_{t=0} D\geo_{-t}(U)=\Phi^- U,\\
&\frac{\diff}{\diff t}\Big\vert_{t=0} D\geo_{t}(V)=\Phi^+ V,
\end{split}
\end{equation}

\begin{proposition}\label{prop:G}
Setting 
\[
G_{t,\omega}(x)=\exp\left( 2\pi \imath \int_0^t \langle \omega, X\rangle\circ \geo_{-a}(x)\diff a  \right),
\]
then for each $t\in \R^+$, $ I\subset \cI_\rho, \omega \in \T^d, \psi\in \cC^r_c(I)$ with $\|\psi\|_{\cC^r(I)}\le C_*$ and $\varphi\in \cC^{1+\alpha}_c(I)$, we have
\[
 \|(\varphi G_{t,\omega})\circ \geo_t\circ \horo_{(\cdot)}\circ \psi \|_{\cC^{\varsigma}(I)}  \le C_\varsigma(t) \|\varphi\|_{\cC^\varsigma(I)}, \quad \varsigma \in \{0,\alpha, 1+\alpha\},
\]
where
\begin{equation}\label{eq:C(t)}
C_\varsigma(t)=\begin{cases}  1 \qquad &\text{if}\qquad  \varsigma=0,\\ 					   
					    2\pi C_*\big(\frac{1 }{\sqrt{-\overline k}}\|\langle \omega, X \rangle \|_{\cC^1} +e^{-\sqrt{-\overline k}\alpha t}\big) \qquad &\text{if} \qquad\varsigma=\alpha,\\
					   C_*+8\pi^2\max\{C_*^3,C_*^4\} B(\overline k)A(X,\omega)\big(1+4e^{-\sqrt{-\overline k}(1+\alpha)t}\big) \qquad &\text{if} \qquad\varsigma=1+\alpha.
		\end{cases}
\end{equation}

\end{proposition}

\begin{proof}
The case $\varsigma=0$ is trivial since $ \|G_{t,\omega}\|_{\cC^0}=1$.
Let us use the symbol $^\prime$ to denote the derivative with respect to $\eta \in I$. We have
\[
\begin{split}
((\varphi G_{t,\omega})\circ & \geo_t\circ \horo_\eta\circ \psi )'\\
%
=&(G_{t,\omega}\circ \geo_t\circ \horo_\eta \circ \psi)'\cdot (\varphi\circ \geo_t\circ \horo_\eta \circ \psi)+ G_{t,\omega}\circ \geo_t\circ \horo_\eta \circ \psi\cdot (\varphi\circ \geo_t\circ \horo_\eta \circ \psi)'\\
%
=&2\pi \imath\psi'\left(  \int_0^t (D\geo_{t-a}(U) \langle \omega, X\rangle)\circ \geo_{t-a}\circ \horo_\eta\circ \psi \diff a   \right) (G_{t,\omega}\cdot\varphi)\circ \geo_t\circ \horo_\eta \circ \psi\\
%
&+G_{t,\omega}\circ \geo_t\circ \horo_\eta \circ \psi\cdot \psi' (D\geo_t(U)\varphi)\circ \geo_t\circ \horo_\eta \circ \psi.
\end{split}
\]
Recalling that $D\geo_tU=J_t U$ and $J_{-t}=D\geo_{-t}(U)$, we note that
\begin{equation}\label{eq:great}
(\varphi\circ \geo_t\circ \horo_\eta \circ \psi)'=\psi' (D\geo_t(U)\varphi)\circ \geo_t\circ \horo_\eta \circ \psi= \psi' J_t (U\varphi)\circ \geo_t\circ \horo_\eta \circ \psi.
\end{equation}
Therefore, since $\|\psi\|_{\cC^r}\le C_*$ and $\|G_{t,\omega}\|_{\cC^0}=1$, we have
\[
\begin{split}
|((\varphi G_{t,\omega})\circ \geo_t\circ \horo_\eta\circ \psi )'| 
&\le 2\pi C_*\|\varphi\|_{\cC^0}\|\langle \omega, X \rangle \|_{\cC^1}\int_0^t |J_{t-a} \circ \horo_\eta \circ \psi |\diff  a+C_*|J_t\circ \horo_\eta|\|\varphi\|_{\cC^1}.
\end{split}
\]
By Lemma \ref{lemma:Jacobi}, $|J_t(x)|\le e^{-\sqrt{-\overline k}t}$ for each $x\in M$, and
\begin{equation}\label{eq:intJ}
\left|\int_0^t J_{t-a}\circ h_\eta \diff a\right|\le \int_0^t e^{-\sqrt{-\overline k}(t-a)}\diff a \le \dfrac{1}{\sqrt{-\overline k}},
\end{equation}
therefore,
\begin{equation}\label{eq:C1G}
\begin{split}
\|(\varphi G_{t,\omega})\circ \geo_t\circ \horo_\eta\circ \psi \|_{\cC^1} 
&\le \frac{2\pi C_*}{\sqrt{-\overline k}}\|\langle \omega, X \rangle \|_{\cC^1} \|\varphi\|_{\cC^0}+C_*e^{-\sqrt{-\overline k}t}\|\varphi\|_{\cC^1}\\
&\le C_*\left(\frac{2\pi  }{\sqrt{-\overline k}}\|\langle \omega, X \rangle \|_{\cC^1} +e^{-\sqrt{-\overline k}t}\right)\|\varphi\|_{\cC^1}.
\end{split}
\end{equation}
For $\varsigma=\alpha$ it is enough to observe that, for each $\eta,\tilde \eta \in (0,\rho)$,
\begin{equation}\label{eq:Holder}
\begin{split}
|(\varphi \circ \geo_t\circ \horo_{(\cdot)}\circ \psi)(\tilde\eta)- (\varphi \circ \geo_t\circ \horo_{(\cdot)}\circ \psi_{\tilde\eta})(\eta)| 
&\le d_{I}(\geo_t\circ \horo_\eta\circ \psi,\geo_t\circ \horo_{\tilde \eta}\circ \psi )^{\alpha} \|\varphi\|_{\cC^\alpha}  \\
&\le C_* e^{-\sqrt{-\overline k} \alpha t}|\eta-\tilde \eta|^{\alpha}  \|\varphi\|_{\cC^\alpha},
\end{split}
\end{equation}
where $d_I(x,y)$ is the distance of two points $x,y$ along $I$. Hence, since $G_{t,\omega}\in \cC^r$, using the inequality
\begin{equation}\label{eq:Calpha}
\|\varphi G\|_{\cC^\alpha} \le \|\varphi\|_{\cC^0}\|G\|_{\cC^1}+\|\varphi\|_{\cC^\alpha}\|G\|_{\cC^0},
\end{equation}
we have
\[
\begin{split}
\|(\varphi G_{t,\omega})\circ & \geo_t\circ \horo_\eta\circ \psi \|_{\cC^\alpha(I)}\le \\
&\le \|G_{t,\omega}\circ \geo_t\circ \horo_\eta\circ \psi\|_{\cC^1}\|\varphi\|_{\cC^0}+\|G_{t,\omega}\|_{\cC^0}\|\varphi\circ \geo_t\circ \horo_\eta\circ \psi\|_{\cC^{\alpha}}\\
&\le  C_*\left(\frac{2\pi }{\sqrt{-\overline k}}\|\langle \omega, X \rangle \|_{\cC^1} +e^{-\sqrt{-\overline k}t}\right)\|\varphi\|_{\cC^0}+C_* e^{-\sqrt{-\overline k} \alpha t}  \|\varphi\|_{\cC^\alpha}\\
&\le C_*\left(\frac{2\pi }{\sqrt{-\overline k}}\|\langle \omega, X \rangle \|_{\cC^1} + e^{-\sqrt{-\overline k} \alpha t}  \right) \|\varphi\|_{\cC^\alpha}.
\end{split}
\]
It remains the case $\varsigma=1+\alpha$. Using \eqref{eq:Calpha}, we can check that if $G\in \cC^2$ and $\varphi \in \cC^{1+\alpha}$, then 
\begin{equation}\label{eq:C1+}
\|\varphi G\|_{\cC^{1+\alpha}}\le  \|\varphi\|_{\cC^{1+\alpha}}(1+4\|G\|_{\cC^{2}}).
\end{equation}
 Let us thus compute the second derivative of $(G_{t,\omega}\circ \geo_t\circ \horo_\eta \circ \psi)$: using again that $D\geo_tU=J_t U$,
\begin{equation}\label{eq:2deriv}
\begin{split}
|(G_{t,\omega}\circ & \geo_t\circ \horo_\eta \circ \psi)''|\\
%
=&\left|\left( 2\pi\imath (G_{t,\omega}\circ \geo_t\circ \horo_\eta\circ \psi)\psi' \int_0^t J_{t-a}\circ \horo_\eta \cdot U(\langle \omega, X \rangle\circ \geo_{t-a}\circ \horo_\eta\circ \psi )\diff a \right)'\right|\\
%
\le & 2\pi C_* \bigg( \left| \int_0^t (J_{t-a}\circ \horo_\eta)' \cdot U(\langle \omega, X \rangle\circ \geo_{t-a}\circ \horo_\eta\circ \psi )\diff a \right|\\
%
&+ \left| \int_0^t J_{t-a}\circ \horo_\eta \cdot (U(\langle \omega, X \rangle\circ \geo_{t-a}\circ \horo_\eta\circ \psi ))'\diff a \right|\bigg)\\
%
&+2\pi C_*  \left| \int_0^t J_{t-a}\circ \horo_\eta \cdot U(\langle \omega, X \rangle\circ \geo_{t-a}\circ \horo_\eta\circ \psi )\diff a \right|(1+\|G_{t,\omega}\circ \geo_t\circ \horo_\eta\circ \psi\|_{\cC^1}).
\end{split}
\end{equation}
We need to estimate the above terms. By \eqref{eq:Jbound}, \eqref{eq:potential}, \eqref{eq:div} and \eqref{eq:intJ},
\[
\begin{split}
\left| (J_{t-a}\circ {\horo_\eta})' \right| &\le \left| \int_0^t (\Phi^{-} \circ \geo_a\circ \horo_\eta)' \right|\left| J_{t-a}\circ \horo_\eta \right|\le  \dfrac{\|\Phi^{-} \|_{\cC^1}}{\sqrt{-\overline k}}e^{-\sqrt{-\overline k}(t-a)}.
\end{split}
\]
On the other hand
\[
\left| (U(\langle \omega, X \rangle\circ \geo_{t-a}\circ \horo_\eta\circ \psi ))'\right| \le C_*\|\langle \omega, X \rangle\|_{\cC^2}\|J_{t-a}\|_{\cC^0}\le C_*\|\langle \omega, X \rangle\|_{\cC^2}e^{-\sqrt{-\overline k}(t-a)}
\]
and
\[
\int_0^t |U(\langle \omega, X \rangle\circ \geo_{t-a}\circ \horo_\eta\circ \psi )|\diff a \le C_*\|\langle \omega, X \rangle\|_{\cC^1}\int_0^t e^{-\sqrt{-\overline k}(t-a)}\diff a.
\]
Using the above three inequalities into \eqref{eq:2deriv}, and recalling \eqref{eq:C1G}, we obtain
\[
\begin{split}
\|G_{t,\omega}\circ  \geo_t\circ \horo_\eta \circ \psi\|_{\cC^2}&\le 2\pi C_*^2 \dfrac{\|\Phi^{-} \|_{\cC^1}\|\langle \omega, X \rangle\|_{\cC^1}}{2{\sqrt{-\overline k}}} + \frac{2\pi C_*^2\|\langle \omega, X \rangle\|_{\cC^2}}{2\sqrt{-\overline k}}\\
%
&+\frac{2\pi C_*^2\|\langle \omega, X \rangle\|_{\cC^1}}{2\sqrt{-\overline k}}\left(1+C_*\left(\frac{2\pi  }{\sqrt{-\overline k}}\|\langle \omega, X \rangle \|_{\cC^1} +e^{-\sqrt{-\overline k}t}\right)\right).
\end{split}
\]
We have thus obtained
\[
\|G_{t,\omega}\circ  \geo_t\circ \horo_\eta \circ \psi\|_{\cC^2}\le 2\pi^2\max\left\{\frac{C_*^2}{\sqrt{-\overline k}}, \frac{ C_*^3}{-\overline k}\right\}A(X,\omega)\left(4+e^{-\sqrt{-\overline k}t}\right).
\]
Hence, by \eqref{eq:C1+}, we have
\begin{multline*}
\|(\varphi G_{t,\omega})\circ \geo_t\circ \horo_\eta\circ \psi \|_{\cC^{1+\alpha}} \\ \le (1+8\pi^2\max\{C_*^2,C_*^3\}B(\overline k)A(X,\omega)(4+e^{-\sqrt{-\overline k}t}))\|\varphi\circ \geo_t\circ \horo_\eta\circ \psi \|_{\cC^{1+\alpha}}.
\end{multline*}
Finally, by \eqref{eq:great}, and arguing as in \eqref{eq:Holder},
\[
\|(\varphi\circ \geo_t\circ \horo_\eta\circ \psi)'\|_{\cC^{\alpha}}=\| \psi' J_t (U\varphi)\circ \geo_t\circ \horo_\eta \circ \psi\|_{\cC^\alpha}\le C_*e^{-\sqrt{\overline k}t}e^{-\sqrt{-\overline k} \alpha t}\|\varphi'\|_{\cC^\alpha}.
\]
Inserting the above in the previous equation we conclude:
\begin{multline*}
\|(\varphi G_{t,\omega})\circ \geo_t\circ \horo_\eta\circ \psi \|_{\cC^{1+\alpha}} \\ \le  \left(C_*+8\pi^2\max\{C_*^3,C_*^4\} B(\overline k)A(X,\omega)\big(1+4e^{-\sqrt{-\overline k}(1+\alpha)t}\big)\right)\|\varphi\|_{\cC^{1+\alpha}}.
\end{multline*}
The proof is therefore complete.
\end{proof}

\end{document}
